\documentclass[11pt]{article}
\usepackage[letterpaper,margin=0.65in]{geometry}
\usepackage[T1]{fontenc}
\usepackage[default]{sourcesanspro}
\usepackage{amsmath,amssymb,enumitem,fancyhdr,xcolor}
\setlength{\emergencystretch}{1em}
\setlength{\parindent}{0pt}\setlength{\parskip}{6pt}
\definecolor{ink}{HTML}{183642}
\pagestyle{fancy}\fancyhf{}\fancyhead[L]{\small\textcolor{ink}{Week 7 / Take-away games / Adult return-visit guide}}
\fancyfoot[L]{\small Bellingham Math Circle / Week 7 / RV7-FAC-v1}\fancyfoot[R]{\small\thepage}
\renewcommand{\headrulewidth}{0.4pt}\renewcommand{\footrulewidth}{0pt}
\setlist[itemize]{leftmargin=1.2em,itemsep=2pt,topsep=2pt}
\newcommand{\parttitle}[1]{{\large\bfseries\color{ink}#1}\par}
\newcommand{\labelled}[2]{\textbf{#1} #2\par}
\begin{document}
\parttitle{Mathematical destination}
\labelled{A shared resource changes the state.}{Take 1 or 2; taking the last counter wins. Without PASS, the losing positive piles are multiples of 3. With one shared PASS still available, the losing piles are 4,7,10,$\ldots$; piles 1,2,3 are winning. PASS consumes a turn and is gone for both players. The game ends immediately on taking the final counter; there is no pass from an empty pile.}
\labelled{A visible memory changes the state.}{Take 1,2 or 3, except the number on LAST, which is the previous player's removal. With no prior move, the losing piles have remainder 0 on division by 4. With LAST=1, the losing remainders are 0 and 1; with LAST=2, only 0; with LAST=3, 0 and 3. The player unable to move loses; a single counter with LAST=1 is therefore losing.}
\labelled{A row is more than a total.}{Removing exactly two neighbouring counters from a fixed row may split it into runs. Gaps never close. Single rows of lengths 1--12 lose exactly at 1,5,9; an empty row is terminal. Every even-length row is winning: remove the central pair and then copy in the equal halves. Two equal rows are losing by the corresponding reply. One row of four wins, but two rows of two lose despite equal totals. No general period for arbitrary row lengths is asserted.}
\labelled{Discovery and explanation.}{Experiments may suggest a favourable start; a guaranteed strategy must answer every legal opponent move. In the first two games the state includes a visible extra item. In the row game it includes the surviving gaps. Adults have exact small-state checks and induction arguments; children can explain guarantees by a reply rule with objects.}
\labelled{Band map.}{Shared PASS invites grades 2--5 after ordinary take-away play. Student page 2 uses grades 4--5; LAST adds a visible memory and reading of three move cards. A ready third grader can enter with adult support. Adjacent-pair removal has a K--5 physical entry: a partner checks neighbours and keeps the spaces fixed. Counting up to 15 and removing small numbers suffice. No binary arithmetic or formal game theory is needed. An adult may read, move a card or scribe; children choose the moves.}
\labelled{Status and novelty.}{Prepared, unpiloted. These are resource, memory and spatial investigations, distinct from the base packet's new lists of allowed amounts, misere subtraction and ordinary two-pile charts. Choose one for a visit and preserve records for another; the theme has no one-hour ceiling.}
\newpage
\parttitle{Prepare, launch and choose a route}
\labelled{Per pair.}{Up to 24 counters, one shared PASS slip, cards LAST=1,2,3 and a blank LAST card, pencils and scrap paper. The printed page 3 uses 0.80-inch active squares, including a 7.20-inch nine-square row. For row games use fixed square positions or a taped strip so empty places stay visible. Counters must fit the printed positions; at most 0.75-inch pieces on 0.8-inch or larger squares are a sensible design. Printing and actual material fit have not been physically rehearsed.}
\labelled{Rules made visible.}{Put PASS between the players; spending it turns the card face down and leaves it unavailable. In the remembered game, move the correct LAST card beside the pile after every removal; the partner checks it before the next move. In row games leave every removed counter's position empty. Do not collect the survivors into a new pile. The partner's role is to check legality rather than supply the next move.}
\labelled{Launch together, 3--5 minutes.}{Let children arrange counters. Show a legal ordinary removal, then spend PASS without touching the pile. Show the partner the face-down card, which cannot be spent again. Before the LAST investigation, demonstrate a practice move from six with LAST=1: remove two, leaving four and LAST=2. Before row play, remove a neighbouring pair from seven fixed positions and leave the two surviving runs separated. Introduce the new convention with the page that uses it, rather than presenting all three rules at once.}
\labelled{First visit.}{Use 15--25 minutes for one game, swapping who starts and taking time to replay a counterexample. Younger children can remain with short row games; a further 20--30-minute visit can compare PASS states or introduce LAST. The current eleven-child group has three anchored adults at KK11 / 3333 / 445 tables; the upper third player can referee and rotate. A continuation need not start on page 1 if its prerequisite is already familiar.}
\labelled{Keep records light.}{Save a starting state and the move sequence, or photograph the final visible state. A result should name both pile size and card state, or all run lengths. Do not demand simultaneous turn tallies, pile tallies and written explanations. Use a chart only when the children want to compare states.}
\labelled{Lineage and evidence.}{Current Week 7 F07-v4 already has ordinary subtraction, differing allowed moves, eventual periodicity, last-counter-loses and two-pile games. The new tasks use original state recurrences. \emph{Math Circle by the Bay}, preface printed pp. viii--x / PDF pp. 9--11, actually supports manipulatives, explanations, reserve challenges and varying depth. These exact games, timings and fixed-table arrangements are our adaptations. No classroom success is claimed.}
\newpage
\parttitle{Problem 1 / The shared PASS}
\labelled{Checked states.}{On the actual 1--12 lists, circle 1,2,3,5,6,8,9,11,12 with PASS available, and 1,2,4,5,7,8,10,11 with PASS used. The preferred player changes at 3,4,6,7,9,10,12. With PASS gone, a pile is losing iff its positive size is divisible by 3. With PASS present, losing starts up to 15 are 4,7,10,13. Piles 1 and 2 win by ending the game immediately; at 3, pass to the opponent and then answer their removal with the remaining counters.}
\labelled{A concrete strategy.}{At an available-card pile congruent to 2 modulo 3 and at least 5, take 1 to leave an available-card loser. At a multiple of 3, either pass to an ordinary losing pile or, for sizes at least 6, take 2 to leave an available-card loser. At an available-card loser, whichever move the opponent chooses leads to one of these winning cases; respond accordingly. Once PASS is gone use the familiar make-a-multiple-of-three strategy.}
\labelled{Why the pattern holds.}{Check 1,2,3 directly. For $n\ge4$, an available-card pile $n\equiv1\pmod3$ can take 1 or 2 only to available winners, and can pass only to an ordinary winner because $n$ is not a multiple of 3. Every other available-card size can reach a losing available-card size, except the small terminal cases already checked; a multiple of 3 can also pass to an ordinary loser. Induction makes the reply strategy a guarantee for all pile sizes.}
\labelled{A boundary trap.}{It is false that every available-card pile with remainder 1 is losing: a pile of one can be taken immediately. Losing labels apply to the next player's whole state, not a number considered in isolation. The shared card also differs from each player owning a private pass; that is another game.}
\labelled{Entry questions and hints.}{Let children play the same pile with the card present and absent. Ask what a pass changes even though no counter moved. If needed, put two copies of the pile side by side, one with PASS and one without, and compare one move from each. Do not reveal the shifted pattern as a printed recipe.}
\labelled{Continuation.}{Invent a different allowed-removal list with the same shared PASS and test its states. The recurrence is: an available state wins if some removal reaches an available loser, or passing reaches a no-PASS loser. A rule found for 1-or-2 need not survive a changed removal list. Save the exact rule and tested range before making a general claim.}
\newpage
\parttitle{Problem 2 / What happened last?}
\labelled{Checked classification.}{For a pile size $n$, write its remainder modulo 4. The complete losing-state table is below; a dash means that state is winning. Empty piles have already ended the game. On the printed 1--8 chart, circle 1,2,3,5,6,7 with no previous move; 2,3,6,7 for LAST=1; 1,2,3,5,6,7 for LAST=2; and 1,2,5,6 for LAST=3.}
\begin{center}\begin{tabular}{c|cccc}
Remainder of $n$ & No prior move & LAST=1 & LAST=2 & LAST=3\\\hline
0 & losing & losing & losing & losing\\
1 & -- & losing & -- & --\\
2 & -- & -- & -- & --\\
3 & -- & -- & -- & losing
\end{tabular}\end{center}
\labelled{Why the same pile differs.}{At a pile of one, LAST=1 forbids the only affordable move, so the next player loses; any other LAST permits taking the last counter. At three with LAST=3, taking 1 leaves two with LAST=1, a winner; taking 2 leaves one with LAST=2, a winner. With no prior move at three, take all three and win. The remembered amount has real mathematical effect.}
\labelled{Why the table repeats.}{At remainder 0, every allowed amount $k$ leaves remainder $4-k$ with LAST=$k$, a winning state in the table. At remainder 1 with LAST=1, the legal 2 or 3 leads to winning states. At remainder 3 with LAST=3, legal 1 or 2 likewise leads to winners. From remainder 1 in other columns, taking 1 reaches a loser. From remainder 3 in other columns, taking 3 reaches a loser. From remainder 2, take 1 unless LAST=1; in that case take 2. Both reach a losing state. Small piles are checked directly before applying this induction; unaffordable moves are omitted.}
\labelled{Verification recurrence.}{For any specific state, list affordable $k\in\{1,2,3\}$ with $k\ne\mathrm{LAST}$. A state wins exactly when one move leads to a losing state $(n-k,k)$. This is the independent program's rule; no numerical pattern is assumed by it.}
\labelled{Hint and return work.}{Replay a disputed start with the LAST card visibly beside it. Ask whether a label attached only to the pile size could ever be enough. Later children can design an opponent's difficult start or try a different forbidden-repeat rule; first retain time for concrete play. Record the starting LAST, not merely the starting pile.}
\newpage
\parttitle{Problem 3 / Pairs in a fixed row}
\labelled{Checked small rows.}{A single row with 1--12 counters favours the second player at 1,5,9 and the first player at 2,3,4,6,7,8,10,11,12. The actual pictured cases are A: one run of 4 (first wins), B: two runs of 2 (second wins), C: one run of 5 (second wins), D: one run of 3 (first wins), E: one run of 6 (first wins), F: one run of 9 (second wins). The printed non-task example removes positions 2 and 3 from seven, leaving runs of 1 and 4 without closing the gap. A row of one has no legal pair. Removing the end pair at 7 leaves the losing row of 5; doing so at 11 leaves the losing row of 9. Every even row has a central-pair first move leaving two equal runs.}
\labelled{Copying is a proof.}{If two surviving rows have equal lengths, every pair removed from one has a corresponding available pair in the other. Reply there. After the reply, the surviving runs again occur in identical pairs. The second player can always reply, so the first player is the first unable to move. Starting from one even row, the central removal creates this paired position for the opponent. Gaps must remain gaps for the correspondence to be valid.}
\labelled{Total count is insufficient.}{One run of four is winning by removing the central pair. Two runs of two are losing: the first player removes a complete run and the second removes the other. Two runs of three are also losing by copying. These positions need no formal Nim calculation.}
\labelled{Why five loses.}{An end-pair move leaves a run of three, which wins by leaving one. Either inner-pair move leaves runs of one and two; the next player removes the pair and leaves isolated singles. Thus every legal first move from five gives the next player a win. For nine, up to reflection the first move leaves runs 7; 1 and 6; 2 and 5; or 3 and 4. Each is winning: from 7 remove an end pair to leave 5; from 1+6 use the centre of 6 to leave two equal runs of 2 plus an isolated single; from 2+5 remove the run of 2 to leave 5; from 3+4 remove an end pair of 4 to leave 3+2. In 3+2 exactly two pair moves remain, so it loses. Thus every first move from nine gives the next player a win. The independent check also enumerates all pair positions, never inferring outcomes from total size.}
\labelled{Exact finite method.}{Represent surviving lengths as a sorted tuple. Removing positions $i+1,i+2$ from a run of length $n$ replaces that run by lengths $i$ and $n-i-2$, deleting zero lengths. A tuple with no length at least two loses. Otherwise it wins iff any legal successor loses. The portable check script computes the pictured instances from this rule. This is facilitator bookkeeping; children can keep the actual counters instead.}
\labelled{Questions, hints and continuation.}{Ask whether two starts with the same number of counters can favour different players. If needed, suggest comparing a central move with an end move, then let the child decide which helps. Further work can use two unequal rows or seek losing single-row lengths beyond 12; do not promise a simple repeating pattern. Preserve a position and a reply rule for the next visit.}

\end{document}
